\begin{abstract}
Semantic entropy outperforms token entropy by 0.155 AUROC on TriviaQA --- then loses by 0.066 AUROC on TruthfulQA.
Same model, opposite orderings.
This reversal is not noise: it exposes a task-structure condition that determines when semantic-level uncertainty estimation helps and when it fails.
We evaluate four major uncertainty proxies --- semantic entropy (SE), token entropy (TE), SelfCheckGPT BERTScore (SCG), and verbalized confidence (VC) --- at Llama-2-7B scale under identical conditions and find that SE outperforms TE by +0.155 AUROC on TriviaQA (95\% bootstrap CI on the gap excludes zero; CIs marginally overlap at 0.063), with the gap mechanistically explained by TE aggregating surface-form variation within semantically equivalent paraphrase clusters (mean intra-cluster TE variance = 7.152 nats$^2$, 71$\times$ the significance threshold).
However, this advantage reverses on TruthfulQA (TE = 0.511 $>$ SE = 0.445), where adversarial misconceptions produce deterministic-wrong outputs that SE's paraphrase-filtering mechanism cannot handle.
We additionally show that BERTScore-based SCG fails as an SE substitute on short factual QA (gap = 0.336 vs.\ 0.03 equivalence threshold) and that verbalized confidence at 7B scale produces a near-degenerate confidence distribution (ECE = 0.430).
Together, these findings establish the task-structure condition for SE's superiority: SE should be preferred when incorrect model outputs are paraphrase-diverse, and TE when they are deterministically wrong --- a scope condition unidentified in prior 65B-scale evaluations.
\end{abstract}
